\documentclass{article}
\usepackage{listings}
\usepackage{xcolor}
\usepackage{textcomp}
\usepackage[T1]{fontenc}
\usepackage{libertinus}
\usepackage[portuguese]{babel}
\usepackage[top=3cm, left=3cm, bottom=2cm, right=2cm]{geometry}

% Custom Colors
\definecolor{codegreen}{rgb}{0,0.6,0}
\definecolor{codegray}{rgb}{0.5,0.5,0.5}
\definecolor{codepurple}{rgb}{0.58,0,0.82}
\definecolor{backcolour}{rgb}{0.95,0.95,0.92}

\lstdefinestyle{mystyle}{
    backgroundcolor=\color{backcolour},
    commentstyle=\color{codegreen},
    morecomment=[l]{//},
    keywordstyle=\color{magenta},
    numberstyle=\tiny\color{codegray},
    stringstyle=\color{codepurple},
    basicstyle=\small\ttfamily,
    breaklines=true,
    captionpos=b,
    numbers=left,
    numberstyle=\tiny\color{codegray},
    numbersep=5pt,
    showstringspaces=false,
    upquote=true
}

\begin{document}

\begin{flushleft}
    Universidade Federal do Rio Grande do Sul - Instituto de Matemática e Estatística \\
    MAT01126 - Algoritmos e Estruturas de Dados \\
    Leonardo Kawai Suarez - 00578511 \\
    Prof. Dr. Rudnei Dias da Cunha
\end{flushleft}

\begin{center}
    \large{Tarefa: Produto de Matrizes em Pascal e Python}
\end{center}

Implementamos o produto de matrizes em Pascal e Python utilizando o algoritmo \texttt{matmat.kji}, conforme os programas abaixo. Notamos que a implementação em Pascal é mais verbosa, exigindo \texttt{begin} e \texttt{end} para delimitar blocos de código e a definição explícita de tipos e tamanhos de matrizes, enquanto a implementação em Python é mais concisa e flexível, aproveitando listas aninhadas. Ambas as implementações seguem a mesma lógica de multiplicação de matrizes, iterando sobre os índices das matrizes para calcular o produto. A indexação em Pascal permite começar em 1, enquanto em Python começa em 0, o que é uma diferença importante a ser considerada ao traduzir algoritmos entre essas linguagens. Além disso, a implementação em Python utiliza compreensão de listas para inicializar a matriz resultante com zeros, enquanto em Pascal é necessário um loop explícito para isso. Por fim, em termos de desempenho, Pascal pode ser mais rápido por ser uma linguagem compilada, enquanto Python é interpretada, mas a diferença de desempenho pode ser insignificante para matrizes pequenas.

\begin{lstlisting}[style=mystyle, language=Pascal, caption={Produto de matrizes em Pascal}]
program multiplicacao_matrizes;

const
    m = 3; { linhas de A e linhas de C }
    r = 2; { colunas de A e linhas de B }
    n = 3; { colunas de B e colunas de C }
    
type
    matriz_A = array[1..m, 1..r] of Integer;
    matriz_B = array[1..r, 1..n] of Integer;
    matriz_C = array[1..m, 1..n] of Integer;
    
procedure matmat_kji(const A: matriz_A; const B: matriz_B; var C: matriz_C);
var
    i, j, k: Integer;
begin
    for i := 1 to m do
        for j := 1 to n do 
            C[i,j] := 0;
    for k := 1 to r do
        for j := 1 to n do
            for i := 1 to m do
                C[i,j] := C[i,j] + A[i,k]*B[k,j];
end;

procedure imprimir_matriz(const C: matriz_C);
var
    i, j: Integer;
    
begin
    for i := 1 to m do
    begin
        for j := 1 to n do
            write(C[i,j]:4, ' ');
        writeln;
    end;
end;

var
    A: matriz_A;
    B: matriz_B;
    C: matriz_C;

begin
    A[1,1] := 1; A[1,2] := 2;
    A[2,1] := 3; A[2,2] := 4;
    A[3,1] := 5; A[3,2] := 6;
    
    B[1,1] := 1; B[1,2] := 2; B[1,3] := 3;
    B[2,1] := 4; B[2,2] := 5; B[2,3] := 6;
    
    matmat_kji(A,B,C);
    writeln('C = A*B =');
    imprimir_matriz(C);
end.
\end{lstlisting}

\begin{verbatim}
    C = A*B =
    9   12   15
    19  26   33
    29  40   51
\end{verbatim}

\begin{lstlisting}[style=mystyle, language=Python, caption={Produto de matrizes em Python}]
m = 3 # linhas de A e linhas de C 
r = 2 # colunas de A e linhas de B
n = 3 # colunas de B e colunas de C 

def matmat_kji(A, B):
    # inicializa a matriz C com zeros
    C = [[0 for _ in range(n)] for _ in range(m)]
    
    for k in range(r):
        for j in range(n):
            for i in range(m):
                C[i][j] = C[i][j] + A[i][k] * B[k][j]
    return C 

def imprimir_matriz(C):
    for linha in C:
        print(" ".join(f"{val:4d}" for val in linha))

A = [
    [1,2],
    [3,4],
    [5,6]
]        

B = [
    [1,2,3],
    [4,5,6]
]

C = matmat_kji(A,B)
print("C =")
imprimir_matriz(C)
\end{lstlisting}

\begin{verbatim}
    C =
       9   12   15
      19   26   33
      29   40   51
\end{verbatim}

\end{document}